% !BIB program = biber
% !TeX root = main.tex
% Loaded before tumbook so the root tumuser.sty can use these values.
\newcommand{\ThesisTitle}{Simulation of Energy-Dependent Scintillation-Light Patterns of Narrow Gamma Beams}
\newcommand{\ThesisSubtitle}{Physical Mechanisms and Parametric Modelling}
\newcommand{\ThesisAuthor}{Anonymous Author}
\newcommand{\University}{Technical University of Munich}
\newcommand{\Department}{TUM School of Natural Sciences}
\newcommand{\ChairName}{Max Planck Institute for Extraterrestrial Physics}
\newcommand{\DegreeProgramme}{Physics}
\newcommand{\ThesisDegree}{Bachelor of Science (B.Sc.)}
\newcommand{\ThesisExaminer}{PD Dr. Jochen Greiner}
\newcommand{\SubmissionDate}{20.09.2026}
\newcommand{\SubmissionCity}{Puchheim}

\PassOptionsToPackage{textfont=palatino}{tumfonts}
\documentclass[a4paper,thesis=student,english,oneside,openany,coverpage=false,titlepage=false]{tumbook}
\KOMAoptions{headings=optiontoheadandtoc,listof=totoc}

% Retain the TUM layout, with more readable body text and shorter lines.
\pgfkeys{/TUM/page/layout/insidepage/.append style={
  /TUM/page/inner=30mm,
  /TUM/page/outer=30mm}}
\KOMAoptions{fontsize=12pt}
\setsansfont{TeX Gyre Heros}
\setkomafont{TUM.threeliner}{\sffamily\upshape\color{TUMBlue}}
\setmonofont{TeX Gyre Cursor}[Scale=MatchLowercase]
\setlength{\emergencystretch}{1em}
\usepackage{microtype}
\usepackage{csquotes}


% Mathematics and physical quantities
\usepackage{mathtools}
% The TUM Palatino option already loads newpxmath, including AMS symbols.
\usepackage{bm}
\usepackage{derivative}
\usepackage{siunitx}
\usepackage[version=4]{mhchem}
\sisetup{
  separate-uncertainty = true,
  per-mode = symbol,
  range-phrase = {--},
  range-units = single
}

% Figures, tables, and captions
\usepackage{graphicx}
\usepackage{pdfpages}
\usepackage{pgf}
\usepackage{import}
% Keep PGF image sidecars relative to the figure being imported.
\newcommand{\thesisfigure}[3]{\resizebox{#1}{!}{\import{figures/#2/}{#3.pgf}}}
\graphicspath{{figures/}{tum/}}
\usepackage{booktabs}
\usepackage{tabularx}
\usepackage{longtable}
\usepackage{multirow}
\usepackage{subcaption}
\usepackage{placeins}
\usepackage{float}
\usepackage{needspace}
\usepackage{pdflscape}

% Lists and source code
\usepackage{enumitem}
\usepackage{xcolor}
\usepackage{listings}
\lstset{
  basicstyle=\ttfamily\small,
  breaklines=true,
  frame=single,
  numbers=left,
  numberstyle=\tiny\color{gray}
}

% Bibliography: Biber handles UTF-8 names and modern BibLaTeX features.
\usepackage[
  backend=biber,
  style=authoryear,
  sorting=nyt,
  sortcites=true,
  giveninits=true,
  uniquename=init,
  maxcitenames=2,
  mincitenames=1,
  maxbibnames=99,
  doi=true,
  url=true,
  isbn=false
]{biblatex}
% Prefer DOI links; retain URLs and eprints for entries without a DOI.
\AtEveryBibitem{%
  \iffieldundef{doi}{}{\clearfield{url}\clearfield{urlyear}\clearfield{urlmonth}\clearfield{urlday}\clearfield{eprint}\clearfield{eprinttype}\clearfield{eprintclass}}%
}
% Show only publication years for non-web sources; preserve access dates.
\AtEveryBibitem{%
  \ifentrytype{online}{}{%
    \clearfield{month}\clearfield{day}%
    \clearfield{endmonth}\clearfield{endday}%
  }%
}
% Allow long bibliography URLs to wrap within the text margins.
\usepackage{xurl}

% tumbook already loads hyperref with unobtrusive links.
\hypersetup{
  hidelinks,
  pdfauthor={\ThesisAuthor},
  pdftitle={\ThesisTitle}
}
\usepackage[nameinlink,noabbrev]{cleveref}
\crefname{equation}{equation}{equations}
\crefname{figure}{figure}{figures}
\crefname{table}{table}{tables}

% Project-specific notation belongs here. LaTeX Workshop parses imported files,
% so these commands are available in autocomplete throughout the project.
\newcommand{\geantfour}{\textsc{Geant4}}
\newcommand{\ceriumbromide}{\cebrthree}
\newcommand{\cebrthree}{\ensuremath{\mathrm{CeBr_3}}}
\newcommand{\bgo}{\ensuremath{\mathrm{Bi_4Ge_3O_{12}}}}
\newcommand{\edep}{\ensuremath{E_{\mathrm{dep}}}}
\newcommand{\expectation}[1]{\ensuremath{\mathbb{E}\!\left[#1\right]}}
\newcommand{\vect}[1]{\ensuremath{\bm{#1}}}

% Use one sequential figure and table counter throughout the thesis.
\makeatletter
\@removefromreset{figure}{chapter}
\renewcommand{\thefigure}{\arabic{figure}}
\@removefromreset{table}{chapter}
\renewcommand{\thetable}{\arabic{table}}
\makeatother

% Let scientific figures share pages with their discussion rather than queue
% for float-only pages. Local barriers in the chapters bound topic drift.
\setcounter{topnumber}{3}
\setcounter{bottomnumber}{2}
\setcounter{totalnumber}{4}
\renewcommand{\topfraction}{0.9}
\renewcommand{\bottomfraction}{0.8}
\renewcommand{\textfraction}{0.05}
\renewcommand{\floatpagefraction}{0.75}
\setlength{\textfloatsep}{12pt plus 3pt minus 2pt}
\setlength{\intextsep}{12pt plus 3pt minus 2pt}
\setlength{\floatsep}{10pt plus 2pt minus 2pt}
\raggedbottom
% Exact German declaration wording from the official TUM title-page ZIP.
% Its original unbreakable signature line overlaps with descriptive placeholders;
% two wrapping columns also accommodate longer real names and dates.
\makeatletter
\renewcommand{\@tum@student@disclamer@block@}{{%
  \begin{otherlanguage}{ngerman}
  \raggedright\noindent
  Ich versichere hiermit, dass ich die von mir eingereichte Abschlussarbeit
  selbstständig verfasst und keine anderen als die angegebenen Quellen und
  Hilfsmittel benutzt habe.\par
  \vspace{2cm}
  \noindent
  \begin{minipage}[t]{0.48\textwidth}
    \raggedright\theAuthorCity\@tum@date@separator@\theDateSubmitted
  \end{minipage}\hfill
  \begin{minipage}[t]{0.48\textwidth}
    \raggedright\@author
  \end{minipage}\par
  \end{otherlanguage}
}}
\makeatother


\addbibresource{bibliography/references.bib}

\title{\ThesisTitle}
\subtitle{\ThesisSubtitle}
\author{\ThesisAuthor}
\degree{\ThesisDegree}
\dateSubmitted{\SubmissionDate}
\examiner{\ThesisExaminer}

\begin{document}

\frontmatter
% Official TUM title-page layout, compiled by .latexmkrc.
% Only its first page is used; the German declaration is typeset below.
\includepdf[pages=1,pagecommand={\thispagestyle{empty}}]{build/titlepage/Deckblatt.pdf}
% Activate the configured body margins, as tumbook's maketitle previously did.
\pagestyle{TUM.insidepage}
\thispagestyle{empty}
\null\vfill
\makeatletter
\@tum@student@disclamer@block@
\makeatother
\cleardoublepage


\chapter{Abstract}

A narrow gamma beam can produce different scintillation-light patterns at the readout. At low energies, the optical photons spread broadly across the readout, while at higher energies, a sharp central peak appears.

This thesis uses Geant4 simulations to investigate how these patterns arise in cerium bromide (\ce{CeBr3}) and bismuth germanate (BGO), examining the effects of gamma energy, incidence angle, and beam width. The results identify interaction depth and the subsequent propagation of scintillation light as the main influences on the observed shape. Deposits close to the readout produce the sharp central contribution, although most detected light remains distributed over a wider area. Gamma scattering modifies both the centre and the tails but has a comparatively small effect on the overall pattern.

Compact parametric models reproduce the main features of these simulated distributions, with remaining limitations for inclined beams. Together, the physical interpretation and mathematical descriptions could provide a basis for future studies of beam wavelength reconstruction.

\par\vspace{1.5\baselineskip}
\noindent\begin{minipage}{\textwidth}
  \centering
  \thesisfigure{0.90\textwidth}{introduction}{energy_heatmaps}
  \par\smallskip
  {\small\raggedright
  Simulated \ce{CeBr3} light patterns at the readout for perpendicular point beams at \num{100} and \qty{662}{\kilo\electronvolt}. Maps are independently normalized to unit probability and share a logarithmic colour scale.\par}
\end{minipage}

\tableofcontents

\mainmatter

\chapter{Introduction}

\section{Motivation}

Measurements of gamma radiation provide information about its sources through the energies and arrival directions of the detected photons. Scintillation detectors make this radiation accessible to optical readout by converting deposited energy into light~\parencite{pedrini2005scintillation}. The spatial distribution of scintillation light in the detector contains additional information about the beam, such as its wavelength~\parencite{ovsiannikova2020light}.

The measured distribution depends on the interactions within the crystal and the propagation of the scintillation light to the readout. Understanding this response is therefore necessary when relating detector signals to the radiation that produced them.

\section{Previous work and open questions}
This work builds on the Geant4 simulation of a \ce{CeBr3} scintillator developed by \textcite{ovsiannikova2020light}. That study showed that the spatial light distribution at the readout changes with gamma energy and incidence angle. Its detector setup provides the starting point for the present investigation.

These observations raise the question of how the light pattern arises from gamma interactions, the resulting energy deposition, and the propagation of scintillation light. The present work examines their influence on the observed distribution.

\section{Aim and scope of this thesis}
This thesis investigates how the two-dimensional light distribution at the readout depends on gamma energy, incidence angle, and beam profile, and explores the physical origin of these changes. The material dependence is examined by comparing \ce{CeBr3} with bismuth germanate, \ce{Bi4Ge3O12} (BGO), in the same geometry.

A simple mathematical model is developed to approximate the detected light distributions. Such a description could provide a basis for later inference of beam properties from measured light distributions.

The study uses a fixed crystal geometry, gamma energies of \qtyrange{50}{1000}{\kilo\electronvolt}, and an ideal readout. No experimental validation is performed.

\chapter{Foundations}
\section{Gamma interactions}

Gamma radiation can interact with matter in a variety of ways. The processes and their probabilities depend on the photon energy and on the composition and density of the material. In the independent-atom approximation, the linear interaction coefficient for a process $j$ is $\mu_j(E)=\sum_a n_a\sigma_{j,a}(E)$, where $n_a$ is the number density of atoms of element $a$ and $\sigma_{j,a}$ is the corresponding atomic cross section. The sum $\mu(E)=\sum_j\mu_j(E)$ gives the total interaction coefficient, and $\ell_\gamma=1/\mu$ is the mean free path before an interaction~\parencite{hubbell1995attenuation}.

In a homogeneous material, the probability that a monoenergetic gamma travels a distance $u$ without interacting is
\begin{equation}
  P_0(u)=\exp[-\mu(E)u].
\end{equation}
The probability of at least one interaction is consequently $1-\exp[-\mu(E)u]$.
This law describes removal from the uncollided beam, whether by absorption or scattering, and gives an exponential first-interaction distribution along the incident direction~\parencite{hubbell1995attenuation}.

\subsection{Scattering}

One of the ways a gamma may interact with matter is by scattering, changing its direction and imparting some of its energy to the material at the interaction site.

\subsubsection{Rayleigh scattering}

Rayleigh scattering is coherent scattering from the atomic electron cloud in which the atom's internal state remains unchanged. In the elastic approximation, recoil-energy transfer is neglected and the photon changes direction without energy transfer~\parencite{hubbell1969photon}.

\subsubsection{Compton scattering}

In Compton scattering, a gamma transfers energy and momentum to a single electron and continues in a different direction with lower energy. A useful approximation is to treat the electron as free and initially at rest. In this model, the energy transfer and the photon deflection are linked by energy and momentum conservation. More precisely, the outgoing gamma energy after scattering at an angle $\chi$ is
\begin{equation}
  E'=\frac{E}{1+(E/{m_e c^2})(1-\cos\chi)},
  \qquad E_{\mathrm{tr}}=E-E',
  \label{eq:compton-energy}
\end{equation}
where $E$ is the incident gamma energy, ${m_e c^2}$ the electron rest energy, and $E_{\mathrm{tr}}$ the energy transferred to the recoil electron~\parencite{compton1923}.

The angular distribution is described by the Klein--Nishina formula for scattering from unpolarized free electrons initially at rest. Integrating over azimuth gives the differential cross section in $\chi$:
\begin{equation}
  \frac{d\sigma}{d\chi}=\pi r_e^2\sin\chi{\left(\frac{E'}{E}\right)}^2
  \left(\frac{E'}{E}+\frac{E}{E'}-\sin^2\chi\right),
  \qquad 0\leq\chi\leq\pi.
  \label{eq:klein-nishina-angle}
\end{equation}
Here, $E'$ follows \cref{eq:compton-energy} and $r_e$ is the classical electron radius. The factor $\sin\chi$ accounts for the solid angle associated with each angular interval~\parencite{klein1929}.

This polar-angle cross section applies independently of the incident photon polarization. However, polarization does affect the azimuthal distribution. For linearly polarized photons, the electric-field direction supplies a preferred axis perpendicular to the photon propagation direction. The scattering plane is the plane containing the incident and outgoing photon directions. At a fixed polar angle with $0<\chi<\pi$, scattering is more probable when the scattering plane is perpendicular to the electric-field axis than when it contains that axis. Integration over all azimuthal angles removes this modulation~\parencite{klein1929}.

The free-electron description neglects both atomic binding and the electron's initial motion. Binding modifies the scattering probability, while the initial electron momentum broadens the relation between scattering angle and outgoing photon energy~\parencite{liang2020transport}. To illustrate the relative energy scales, the transfer $E_{\mathrm{tr}}$ can be compared with the binding energy $B$. For example, in a \ce{Bi4Ge3O12} (BGO) scintillator, the $K$-shell binding energies of oxygen, germanium, and bismuth are approximately \qty{0.54}{\kilo\electronvolt}, \qty{11}{\kilo\electronvolt}, and \qty{91}{\kilo\electronvolt}, respectively~\parencite[table 1-1]{williams2009binding}. Requiring $B/E_{\mathrm{tr}}<0.1$ gives transfer scales above approximately \qty{5.4}{\kilo\electronvolt}, \qty{111}{\kilo\electronvolt}, and \qty{910}{\kilo\electronvolt}. This shows that even at several hundred \unit{\kilo\electronvolt} of gamma energy, binding effects may still be relevant, though the exact influence depends on scattering angle and electron momentum~\parencite[eqs.~(7)--(11)]{liang2020transport}. For example, a gamma of \qty{662}{\kilo\electronvolt} scattering by \qty{30}{\degree} from lead, with a $K$-shell binding energy of \qty{88}{\kilo\electronvolt}~\parencite[table 1-1]{williams2009binding}, was shown to have a scattering cross section half that predicted by the free-electron model~\parencite[table I]{wolff1989compton}.

\subsection{Photoelectric absorption}

In photoelectric absorption, the gamma is absorbed by an atom and an electron is ejected from its shell. Neglecting atomic recoil, the electron carries away a kinetic energy of $T_e=E-B$, where $B$ is the shell binding energy. The remaining energy is initially associated with the vacancy in the atomic electron configuration~\parencite{ballabriga2016review}.

The vacancy can be filled by an electron from a higher shell. The released energy can be emitted as a characteristic fluorescence photon or transferred directly to another electron, which is then emitted as an Auger electron~\parencite{ballabriga2016review}. These transitions leave further vacancies and can produce a relaxation cascade~\parencite{soronen2024photoionizationauger}. The energies and relative probabilities of these two processes depend on the element and the shells involved~\parencite{ballabriga2016review,soronen2024photoionizationauger}.

When the photon energy reaches the binding energy of an additional electron shell, a further photoionization channel becomes available, causing a sharp increase in the cross section. This is known as an absorption edge. Different shells have different thresholds, so an element has several absorption edges. Between edges, the cross section generally decreases with increasing photon energy. In a compound, the mass attenuation coefficient can be approximated by a sum of the elemental coefficients weighted by their mass fractions, so the compound retains the edges of its constituents. Close to an edge, the measured cross section can show fine structure that depends on chemical state, phase, and temperature. Such structure is omitted in the independent-atom approximation~\parencite[section 2]{hubbell1995attenuation}.

The photoelectric cross section depends strongly on photon energy and atomic number. Away from absorption edges, a rough empirical scaling of the cross section per atom is $\sigma_{\mathrm{pe}} \propto Z^{4.5}/E^3$ over the range \qtyrange{10}{500}{\kilo\electronvolt}~\parencite[table 2-5 and accompanying text]{hubbell1969photon}.

\subsection{Pair production}

In the field of a nucleus, a gamma can convert into an electron--positron pair, while the nucleus takes up recoil momentum. Energy conservation gives a threshold of approximately $2m_ec^2\approx\qty{1.022}{\mega\electronvolt}$ in the heavy-nucleus limit. Energy above the threshold is shared between the kinetic energies of the produced particles and nuclear recoil~\parencite{betheheitler1934pairs}. At even higher energies, pair production can also occur in the field of an atomic electron~\parencite{wheeler1936triples}.

The positron can then annihilate with an electron into two photons~\parencite[p.~347]{klemperer1934annihilation}. For negligible initial kinetic energy, energy and momentum conservation give photons of approximately \qty{511}{\kilo\electronvolt} each, travelling in opposite directions.

\subsection{Energy dependence of interaction probabilities}

\begin{figure}[H]
  \centering
  \thesisfigure{\textwidth}{foundations}{bgo_attenuation}
  \caption{Mass attenuation coefficients of BGO on a double-logarithmic scale. Lines connect the tabulated values. Pair production combines the nuclear- and electron-field contributions. Absorption-edge fine structure is not resolved. The total is the sum of all plotted process coefficients~\parencite{nistxcom}.}
  \label{fig:foundations-attenuation}
\end{figure}

The relative importance of the different gamma interactions is illustrated by the mass attenuation coefficients of BGO in \cref{fig:foundations-attenuation}. The elemental contributions of oxygen, germanium, and bismuth are combined according to their mass fractions. Dividing the linear coefficient by the density gives $\mu_j/\rho$, which expresses the interaction probability per unit mass thickness.

\Cref{fig:foundations-attenuation} shows photoelectric absorption dominating at low energies, with sharp changes at the constituent absorption edges. Compton scattering decreases more slowly with energy and exceeds the photoelectric contribution at approximately \qty{500}{\kilo\electronvolt}. Rayleigh scattering contributes little to attenuation compared with other processes and falls below Compton scattering between \qty{100}{\kilo\electronvolt} and \qty{150}{\kilo\electronvolt} in BGO. Pair production becomes possible above its threshold, although it does not immediately become the dominant process: at \qty{2}{\mega\electronvolt}, the Compton contribution is still larger. By \qty{10}{\mega\electronvolt}, pair production exceeds Compton scattering~\parencite{nistxcom}.

\section{Secondary-particle transport and energy deposition}
\subsection{Difference between energy transfer and deposited energy}

After a gamma transfers energy to an electron through scattering or absorption, the electron carries kinetic energy and loses it through further interactions, including ionization and excitation. Some energy can also return to the photon field through bremsstrahlung or fluorescence~\parencite{hubbell1995energy}. The distribution of deposited energy therefore depends on the transport behaviour of these secondary particles as well as on the positions of the primary gamma interactions.

In a finite crystal, not all energy transferred by a gamma is necessarily deposited in the crystal. A photoelectron or a fluorescence photon can escape the crystal~\parencite{storm2020photoelectronescape, ballabriga2016review}. This distinction is particularly relevant near a surface, where even a short secondary-particle path can carry energy outside the material.

\subsection{Electron stopping and spatial spreading}

The mean rate at which an electron loses kinetic energy along its path is described by the stopping power $S(T)=-\langle\mathrm{d}T/\mathrm{d}u\rangle$, where $T$ is the electron kinetic energy and $u$ the distance travelled. Its collision contribution, $S_{\mathrm{col}}$, arises mainly from inelastic interactions with the electrons of the material. These can excite the target, leaving it in a higher-energy state, or ionize it, releasing another electron. The latter may carry part of the transferred energy away from the original electron track~\parencite[pp. 1--2]{berger1983ranges}.

An electron also loses energy by emitting bremsstrahlung as it is deflected in the electric fields of nuclei and atomic electrons. The corresponding stopping power is $S_{\mathrm{rad}}$~\parencite{berger1983ranges}. The bremsstrahlung photon can propagate independently of the electron and be absorbed elsewhere, extending the spatial redistribution of energy beyond the charged-particle track.

The total stopping power is $S=S_{\mathrm{col}}+S_{\mathrm{rad}}$. The relative size of these terms depends on both electron energy and material. Collision losses are important as the electron slows down, whereas radiative losses become increasingly important at high energies. The nuclear contribution to bremsstrahlung also depends strongly on atomic number, making it more important in heavy elements. The collision contribution depends on the number of target electrons and their binding, commonly represented through the mean excitation energy of the medium~\parencite{berger1983ranges}.

While losing energy, an electron also changes direction through scattering in the material~\parencite[eqs.~(17), (20)]{liang2020transport}.
Many small deflections accumulate along its path, while less frequent large-angle collisions can produce substantial changes in direction. The angular distribution of this multiple scattering depends on the electron energy and on the scattering properties and number density of the atoms. The nuclear charge and its screening by the atomic electrons enter the single-scattering cross section, from which the accumulated angular distribution is obtained~\parencite{lewis1950multiple}.

A compact measure of this change in direction is the mean scattering cosine. For independent elastic collisions in a homogeneous medium, over a path length $u$ for which the electron energy can be treated as constant, the first angular moment of the scattering equation gives
\begin{equation}
  \langle\cos\theta\rangle=\exp\!\left(-\frac{u}{\lambda_{\mathrm{tr}}}\right),
  \qquad
  \frac{1}{\lambda_{\mathrm{tr}}}
  =\sum_i n_i\int_{4\pi}(1-\cos\chi)
  \frac{\mathrm{d}\sigma_i}{\mathrm{d}\Omega}\,\mathrm{d}\Omega.
\end{equation}
Here, $\theta$ is the angle relative to the initial electron direction, $\chi$ the deflection in a single collision, and $\mathrm{d}\sigma_i/\mathrm{d}\Omega$ the differential elastic-scattering cross section for atomic species $i$ with number density $n_i$. The transport mean free path $\lambda_{\mathrm{tr}}$ characterizes the loss of directional alignment. The factor $1-\cos\chi$ gives larger deflections greater weight than small ones. As $u$ increases, the mean cosine decreases, indicating reduced alignment with the initial direction. For a slowing electron, the energy dependence is included by replacing $u/\lambda_{\mathrm{tr}}$ with $\int_0^u\mathrm{d}u'/\lambda_{\mathrm{tr}}[T(u')]$, treating the energy loss as continuous~\parencite[sections II--III]{lewis1950multiple}.

The energy lost over a given path length also fluctuates because individual inelastic collisions transfer different amounts of energy. These fluctuations are called energy-loss straggling~\parencite[sections II.A.3 and II.C]{seltzer1974transmission}. Electrons with identical initial energies can consequently retain different energies after travelling the same distance. Together with the random angular deflections, these fluctuations produce different tracks and stopping positions even for identical starting conditions.

\subsection{Relevant transport scales}

The transport range of a secondary electron can be characterized by the continuous-slowing-down approximation (CSDA). Here, the electron is assumed to lose energy continuously at the mean rate given by the total stopping power. Its path length from an initial kinetic energy $T_0$ to rest is then
\begin{equation}
  R_{\mathrm{CSDA}}(T_0)=\int_0^{T_0}\frac{\mathrm{d}T}{S_{\mathrm{col}}(T)+S_{\mathrm{rad}}(T)}.
\end{equation}
A CSDA range approximates the total distance travelled along a track~\parencite[section 10]{berger1983ranges}. Multiple scattering changes the electron trajectory, so the typical displacement from the starting point is shorter than the total path length.

BGO, used in the previous example, has a density of $\rho=\qty{7.1}{\gram\per\cubic\centi\metre}$ and mass CSDA ranges of \qty{0.027}{\gram\per\square\centi\metre} and \qty{0.71}{\gram\per\square\centi\metre} at electron kinetic energies of \qty{100}{\kilo\electronvolt} and \qty{1}{\mega\electronvolt}~\parencite{nistestar}. Dividing by $\rho$ gives path lengths of approximately \qty{0.037}{\milli\metre} and \qty{0.99}{\milli\metre}, respectively.

In comparison, the mass-weighted sum of the elemental attenuation coefficients in BGO gives $\mu/\rho\approx\qty{4.0}{\square\centi\metre\per\gram}$ at \qty{100}{\kilo\electronvolt} and \qty{0.068}{\square\centi\metre\per\gram} at \qty{1}{\mega\electronvolt}~\parencite{nistxcom}. With $\ell_\gamma=[\rho(\mu/\rho)]^{-1}$, these correspond to mean free paths of approximately \qty{0.35}{\milli\metre} and \qty{20.6}{\milli\metre}, respectively. Therefore, the gamma mean free path is substantially longer than the electron CSDA path length.

Fluorescence photons also carry energy away from the primary interaction site. In BGO, the bismuth $K\alpha_1$ line lies at approximately \qty{77}{\kilo\electronvolt}~\parencite[table 1-2]{kortright2009xray}. This fluorescence-photon energy is therefore comparable to that of the \qty{100}{\kilo\electronvolt} gamma example. Notably, at \qty{80}{\kilo\electronvolt}, the same mass attenuation mixture calculation gives $\mu/\rho\approx\qty{1.9}{\square\centi\metre\per\gram}$ and a mean free path of approximately \qty{0.74}{\milli\metre}~\parencite{nistestar,nistxcom}. This is longer than the mean free path at \qty{100}{\kilo\electronvolt} because the bismuth $K$ absorption edge lies between the two energies, at \qty{91}{\kilo\electronvolt}~\parencite{nistxcom}. Absorption from that shell is possible at \qty{100}{\kilo\electronvolt} but not at \qty{80}{\kilo\electronvolt}, increasing attenuation above the edge.

Unlike characteristic fluorescence lines, bremsstrahlung forms a continuous spectrum whose shape depends on the electron energy and target material. As an experimental example, the energy-weighted spectrum measured from a tungsten target at an accelerating voltage of \qty{100}{\kilo\volt} shows a pronounced maximum near \qty{30}{\kilo\electronvolt}~\parencite{ehrlich1955bremsstrahlung}. In this example, the energy at the bremsstrahlung maximum is of the same order of magnitude as the fluorescence-photon energy discussed above.

\section{Scintillation-light production}

As secondary particles lose energy in a scintillator, part of the deposited energy is emitted as optical photons. The emission occurs at luminescent centres: sites in the crystal with electronic states that can receive excitation energy and release it as light. An excited centre emits a photon when it returns to a lower electronic state through a radiative transition~\parencite{pedrini2005scintillation}.

\subsection{Microscopic origin of scintillation}

In an inorganic insulating scintillator, the microscopic description of scintillation starts with the electronic states of the crystal. The occupied valence states and the conduction states are separated by a band gap $E_g$. Energy deposited by a charged particle can elevate a valence electron into the conduction band, leaving a hole in the valence band~\parencite{pedrini2005scintillation}. Both the initial energetic electron and newly excited conduction electrons can produce further electron--hole pairs if they retain enough energy~\parencites[pp.~38, 44]{shockley1961junctions}.

After this sequence of ionizations and excitations, electrons and holes can still have excess energy relative to the band edges. They exchange energy with lattice vibrations, described as phonons, and approach thermal equilibrium with the crystal. This transfers part of the electronic excitation energy to the lattice before optical emission~\parencite{pedrini2005scintillation}. During this process, an electron travels a distance known as the thermalization distance. Calculations for several halide scintillators give typical thermalization distances for unrecombined electrons of \qtyrange{25}{50}{\nano\metre}, with some electrons travelling a few hundred nanometres~\parencite{wang2012thermalization}.

An electron and a hole can form a bound state, called an exciton. The pair changes the forces on nearby ions and can shift them from their equilibrium positions. If this rearrangement lowers the energy sufficiently, it binds the pair to that region, forming a self-trapped exciton. Energy can migrate as the pair and its surrounding distortion move between neighbouring sites, allowing light to be emitted away from the formation site~\parencite{prange2013excitations}. As an example, diffusion-model fits to the luminescence decay of \ce{SrF2} nanoparticles estimated mean self-trapped-exciton migration lengths of \qtyrange{13}{17}{\nano\metre}~\parencite{chylii2018diffusion}. At high densities, interactions between excitons can remove energy from the light-producing channel. This effect is known as density-dependent quenching~\parencite{prange2013excitations}.

Carrier trapping at defects can also alter the route to light emission.
A defect can introduce a localized electronic state in which an electron or hole can be bound. Capture into such a state traps the carrier near the defect instead of allowing it to move freely through the crystal, changing the routes available for recombination~\parencite{vanginhoven2013defects}.

In BGO, measurements of the light intensity and decay time show that nonradiative decay becomes significant above approximately $\qty{250}{\kelvin}$. The temperature dependence is attributed to thermally activated nonradiative processes competing with radiative decay~\parencite{gironnet2008bgo}.

The radiative channel terminates with light emission from a luminescent centre. This can occur either when an electron--hole pair recombine, acting as an intrinsic luminescent centre or when the excitation energy is transferred to another luminescent centre, which emits through a transition between its own electronic states. Recombination does not always produce a photon: a nonradiative route transfers energy to the lattice instead~\parencite{pedrini2005scintillation}.

\subsection{Conversion of deposited energy to photon number}

Since some of the energy is diverted to nonradiative processes, it is necessary to consider the scintillation yield, which expresses the mean number of emitted photons per unit deposited energy. It depends on the efficiency of creating electronic excitations, their transfer to emitting centres, and the radiative efficiency of those centres~\parencite{pedrini2005scintillation}.

For a proportional scintillator, the mean emitted photon number is $\langle N_\gamma\rangle=Y E_{\mathrm{dep}}$, with constant yield $Y$. A nonproportional response instead has a light yield per unit deposited energy that varies with the energy of the exciting particle. In practice, fully proportional scintillators do not exist. Nevertheless, a constant yield can be a useful approximation whose accuracy depends on the material, energy range, and external conditions. For example, room-temperature measurements of BGO reported about \qty{14}{\percent} variation in light yield per unit energy over gamma energies of \qtyrange{32}{1274}{\kilo\electronvolt}, normalized to the yield at \qty{662}{\kilo\electronvolt}~\parencite{sreebunpeng2012comparison}.

One cause of nonproportionality is density-dependent quenching. Since the excitation density varies along an electron track and with electron energy, the light emitted per unit deposited energy can vary~\parencite[section III]{payne2009nonproportionality}.

\subsection{Emission spectrum}

Scintillation photons have a material-dependent spectrum whose shape reflects the radiative transitions responsible for luminescence. These can involve centres intrinsic to the host crystal or ions introduced as activators. Their electronic states and coupling to the surrounding lattice differ, so different scintillators have different emission bands~\parencite{prange2013excitations}.

\section{Optical transport to the readout}

Even for an isotropic scintillation source, the light arriving at a detector plane is not uniform. An area element $\mathrm{d}A$ at distance $R$ from the source subtends the solid angle $\mathrm{d}\Omega=\cos\alpha\,\mathrm{d}A/R^2$, where $\alpha$ is the angle between the incoming ray and the plane normal. For a point source a perpendicular distance $d$ from the plane, $R^2=d^2+r^2$ and $\cos\alpha=d/R$, where $r$ is the distance from the projected source position. The direct photon-arrival probability per unit area is therefore $\mathrm{d}P/\mathrm{d}A=d/[4\pi(d^2+r^2)^{3/2}]$, before absorption or reflections.

During propagation, an optical photon may be absorbed. As for gamma radiation, in a homogeneous medium with absorption length $\Lambda(\lambda)$, the survival probability after an optical path of length $u$ is $\exp[-u/\Lambda(\lambda)]$. Here, $\Lambda$ is the mean distance travelled before absorption and in general depends on wavelength~\parencite[fig.~8]{miller2025material}.

Optical photons may travel distances much greater than the size of the scintillator because of reflections at its boundaries, including total internal reflection. The latter follows from Snell's law, $n_1\sin\alpha_1=n_2\sin\alpha_2$, where $n_1$ is the refractive index inside the scintillator and $n_2$ that of the adjacent material. For $n_1>n_2$, incidence above $\alpha_c=\arcsin(n_2/n_1)$ gives total internal reflection.

Not all scintillation photons that reach the readout produce a signal. The photon detection efficiency is the probability that an arriving photon produces a detectable response and generally depends on wavelength~\parencite{eckert2010characterisation}. A spatial photon-arrival distribution must therefore be combined with the readout's wavelength-dependent detection efficiency to determine the expected detected distribution.

\section{Photodetection with silicon photomultipliers}

A commonly used detector of scintillation light is the silicon photomultiplier (SiPM). It consists of many small photodiodes, called microcells. Each photodiode contains a junction between p-type silicon, where holes are the majority mobile carriers, and n-type silicon, where electrons are the majority mobile carriers. Photon absorption can create an electron--hole pair. An applied voltage separates the carriers and drives them towards the electrical contacts. Within the high-field multiplication region, a carrier can gain enough energy to create further electron--hole pairs by impact ionization, producing an avalanche~\parencite{piatek2014sipm}.

The breakdown voltage $V_{\mathrm{br}}$ marks the onset of self-sustaining avalanche multiplication. Operating a microcell at a voltage $V>V_{\mathrm{br}}$ is called Geiger mode: a single carrier can then trigger an avalanche that must be stopped externally~\parencite{piatek2014sipm}.
In a passively quenched SiPM, a series resistor reduces the voltage during the discharge until the avalanche stops. The cell then recharges towards its initial operating voltage. In a simple circuit approximation, the recovery time is $\tau\approx R_q C_{\mathrm{cell}}$, where $R_q$ is the quenching resistance and $C_{\mathrm{cell}}$ the effective pixel capacitance. A further avalanche during recovery produces less charge because the overvoltage $V_{\mathrm{bias}}-V_{\mathrm{br}}$ has not yet been fully restored~\parencite{eckert2012response}.

The photon detection efficiency (PDE) includes several probabilities. A useful factorization is
\begin{equation}
  \mathrm{PDE}(\lambda,V)=f_{\mathrm{fill}}\,\eta(\lambda)\,P_{\mathrm{av}}(\lambda,V),
\end{equation}
where $f_{\mathrm{fill}}$ is the fraction of the sensor area occupied by photosensitive cells, $\eta$ the probability of producing a charge carrier available for detection there, and $P_{\mathrm{av}}$ its probability of initiating an avalanche. Here, $\lambda$ is the photon wavelength and $V$ the applied bias voltage. The PDE therefore varies with photon wavelength~\parencite{eckert2010characterisation}.

At high light levels, several photons can reach the same microcell before it recovers. The output then ceases to be proportional to photon number and approaches saturation~\parencite{eckert2012response}.

SiPM noise includes both independently initiated avalanches and additional avalanches caused by an earlier one. Thermally generated carriers can initiate dark counts without an incident photon. Any avalanche, whether initiated by light or by a dark event, can in turn produce correlated noise: photons generated during the avalanche can trigger a neighbouring cell, causing optical crosstalk, while trapped carriers released later can produce afterpulses~\parencite{eckert2010characterisation}.

% chktex-file 29
\chapter{Simulations}
\label{chap:simulations}

\section{Simulation framework}

The simulations use Geant4 11.4.1, a Monte Carlo particle-transport simulation toolkit~\parencite[Geant4 Scope of Application]{geant4introduction}. Primary gammas are generated with a specified position, direction, and energy and tracked until they are absorbed or leave the simulated volume. The starting configuration is FTFP\_BERT, a predefined collection of particle-interaction models supplied by Geant4~\parencite{geantftfpbert}. Its default electromagnetic models are replaced by polarized Livermore electromagnetic physics, and optical physics is added.
Photoelectric absorption, Compton and Rayleigh scattering, secondary-particle transport, atomic relaxation, scintillation, Auger emission and cascades, and optical transport are computed. The Compton model used here includes shell binding energies and Doppler broadening from bound-electron momentum~\parencite{geant4polarizedcomptonsource}. For each energy-deposition step, the mean scintillation-photon count is $\langle N_\gamma\rangle=Y E_{\mathrm{dep}}$, where $Y$ is the material's constant yield and $E_{\mathrm{dep}}$ the energy deposited in that step. Counts are sampled from a Poisson distribution for means of at most ten photons and from a Gaussian distribution above this threshold~\parencite[Scintillation]{geantopticalprocesses}. No quenching or particle-dependent nonproportionality is simulated. The source code for the simulation application is available on GitHub.\footnote{\url{https://github.com/Lantaan/scintillator-sim}}

Two types of simulation samples are used. Samples used to study the light distribution at the readout contain \num{20000} incident gammas and include scintillation and optical-photon transport processes. Separate samples of \num{100000} incident gammas record the gamma interaction sequences and energy deposited by the gammas and their secondary particles. These larger samples omit scintillation-photon generation to reduce computation time, while gamma scattering and secondary-particle transport remain enabled.

\section{Detector model}

The crystal geometry, reflector arrangement, ideal photon-absorbing readout, and \ce{CeBr3} material and surface models follow the predecessor setup of \textcite[pp.~21--23]{ovsiannikova2020light}. The crystal is enclosed by a reflector and aluminium housing on five faces, with the readout on the remaining face.

\subsection{Reference geometry and irradiation}

The \qtyproduct{51 x 51 x 10}{\milli\metre} crystal is centred at the origin and has thickness $T=\qty{10}{\milli\metre}$. Its entrance is at $z=-T/2=\qty{-5}{\milli\metre}$, and its readout is at $z=T/2=\qty{5}{\milli\metre}$. The transverse coordinates are $x,y$, and the depth measured from the entrance along $z$ is $s=z+T/2$.

An unpolarized monoenergetic point beam starts at $(0,0,-\qty{5}{\milli\metre})$ and travels along $+z$. Inclination by $\theta$ changes the direction to $(\sin\theta,0,\cos\theta)$ without moving the source. The direct readout intersection is $x_{\mathrm{exit}}=T\tan\theta$, $y_{\mathrm{exit}}=0$.

\subsection{Reflector, housing, and optical readout}

The enclosure has a \qty{1}{\milli\metre} PTFE reflector and \qty{0.5}{\milli\metre} aluminium housing on the five non-readout faces. Reflection uses a dielectric--dielectric surface with a reflectivity of 0.97, a zero roughness parameter, and a refractive index matched to the crystal.

A continuous silicon readout volume with a thickness of \qty{1.5}{\milli\metre} covers the remaining face. Its refractive index matches that of the crystal, and its optical absorption length is \qty{0.001}{\milli\metre}, so the absorption coordinates closely match the photon arrival positions. This ideal readout has no physical pixel segmentation, wavelength-dependent detection efficiency, or electronic response.

Unless stated otherwise, comparisons at a single gamma energy use \qty{662}{\kilo\electronvolt}, corresponding approximately to the \ce{^{137}Cs} line~\parencite[section 2.2]{lnhb2023cs137}. Selected-energy comparisons use \qty{100}{\kilo\electronvolt}, \qty{662}{\kilo\electronvolt}, and \qty{1000}{\kilo\electronvolt} to represent the low- and high-energy responses.
\subsection{Inherited \ce{CeBr3} model}

The \ce{CeBr3} model is configured with a density of \qty{5.20}{\gram\per\cubic\centi\metre}, a yield of 60 photons per \unit{\kilo\electronvolt} deposited, and a constant refractive index of 1.9. The emission maximum is near \qty{370}{\nano\metre}. Its optical absorption length decreases from \qty{155}{\milli\metre} at \qty{2.86}{\electronvolt} to \qty{105}{\milli\metre} at \qty{3.60}{\electronvolt}. This inherited curve uses \ce{LaBr3} attenuation data, including absorption and scattering, as a surrogate, following the predecessor model~\parencite[p.~22, table 2.1]{ovsiannikova2020light}.

\chapter{Results}
% chktex-file 29
\section{Characterization of the \ce{CeBr3} optical-photon distribution}

This section establishes the main energy- and geometry-dependent features of the detector-plane optical distribution before their physical origin and parametrization are examined in subsequent sections.

\subsection{Dependence on incident gamma energy}
\label{sec:cebr-energy-scan}

\begin{figure}[H]
  \centering
  \captionsetup{font=small}
  \thesisfigure{0.90\textwidth}{chapter6_1}{figure_6_1_energy_maps_profiles}
  \caption{Unit-normalized detector-plane optical-photon probability densities for \ce{CeBr3} under perpendicular point-beam irradiation at \qty{100}{\kilo\electronvolt}, \qty{662}{\kilo\electronvolt}, and \qty{1000}{\kilo\electronvolt} (top), with the corresponding radial probability densities (bottom). The maps contain exact \qtyproduct{1 x 1}{\milli\metre} cells over $\qty{-25.5}{\milli\metre}\leq x,y\leq\qty{25.5}{\milli\metre}$ and share a logarithmic colour scale. Radial bins have widths of \qty{0.05}{\milli\metre} for $\qty{0}{\milli\metre}\leq r<\qty{1}{\milli\metre}$, \qty{0.25}{\milli\metre} for $\qty{1}{\milli\metre}\leq r<\qty{5}{\milli\metre}$, \qty{1}{\milli\metre} for $\qty{5}{\milli\metre}\leq r<\qty{25}{\milli\metre}$, and \qty{0.5}{\milli\metre} in the final interval to \qty{25.5}{\milli\metre}. Lines connect radial-bin centres.}
  \label{fig:energy-maps}
\end{figure}

The dependence of the optical distribution on the incident gamma energy is first examined with a perpendicular point beam. The \qtyproduct{51 x 51 x 10}{\milli\metre} \ce{CeBr3} crystal is irradiated with \num{20000} gammas at \qty{100}{\kilo\electronvolt}, \qty{662}{\kilo\electronvolt}, and \qty{1000}{\kilo\electronvolt}. Every detector-plane optical map is normalized to unit detected-photon probability, so only its shape is compared. Radial profiles are constructed from detected-photon positions using $r=\sqrt{x^2+y^2}$. The count in each bin is divided by the full-map photon count and the annular area. The radial limit of \qty{25.5}{\milli\metre} excludes incomplete annuli, while photons in the detector corners remain in the normalization.


\Cref{fig:energy-maps} shows a broad, smooth distribution at low gamma energy and a sharp central peak at higher energy, consistent with the qualitative energy dependence reported by \textcite{ovsiannikova2020light}. The radial full width at half maximum (FWHM) is \qty{8.9}{\milli\metre}, \qty{0.23}{\milli\metre}, and \qty{0.22}{\milli\metre}, respectively. Across these three cases, $P(r>\qty{5}{\milli\metre})$ decreases from 0.83 to 0.76 and 0.75, confirming the visual impression that the distribution steepens with increasing energy and becomes nearly identical at \qty{662}{\kilo\electronvolt} and \qty{1000}{\kilo\electronvolt}.



\begin{figure}[H]
  \centering
  \thesisfigure{\textwidth}{chapter6_1}{figure_6_1_depth_conditioning}
  \caption{Depth-conditioned optical distributions for \ce{CeBr3} at normal incidence. The upper row compares the full \qty{100}{\kilo\electronvolt} optical map with \qty{662}{\kilo\electronvolt} maps conditioned on a first interaction in $\qty{-5}{\milli\metre}\leq z<\qty{-4}{\milli\metre}$ and $\qty{4}{\milli\metre}\leq z\leq\qty{5}{\milli\metre}$. The maps are normalized independently, use exact \qtyproduct{1 x 1}{\milli\metre} cells, and share a logarithmic scale. The lower panel shows radial densities for these three samples and for the unconditioned \qty{662}{\kilo\electronvolt} sample. All four profiles use the radial-bin sizes specified in \cref{fig:energy-maps}. Counts are divided by the annular area and by the total detected-photon count of the corresponding sample, including photons in the detector corners. Lines connect radial-bin centres. For the full \qty{100}{\kilo\electronvolt} and entrance-side \qty{662}{\kilo\electronvolt} maps, respectively, the FWHM values are \qty{8.9}{\milli\metre} and \qty{14}{\milli\metre}, $P(r>\qty{5}{\milli\metre})$ is 0.83 and 0.76, and the Pearson correlation between corresponding pixel probabilities is 0.998.}
  \label{fig:depth-conditioning}
\end{figure}

One explanation for the energy dependence is the gamma interaction-depth distribution. Low-energy gammas interact mainly near the entrance face, whereas higher-energy gammas more often reach the detector-side part of the crystal. The deposited-energy field describes where energy is deposited in the crystal, including the effects of secondary-particle transport. The fraction of the total in-crystal deposited energy lying in the detector-side volume $\qty{4}{\milli\metre}<z\leq\qty{5}{\milli\metre}$ rises from \qty{0.15}{\percent} at \qty{100}{\kilo\electronvolt} to approximately \qty{9}{\percent} at both \qty{662}{\kilo\electronvolt} and \qty{1000}{\kilo\electronvolt}. Scintillation light produced closer to the detection plane travels a shorter optical path and is therefore expected to spread less transversely, possibly explaining the change in the optical distribution.

To test this explanation, two \qty{662}{\kilo\electronvolt} optical maps are constructed by selecting events according to their first-interaction depth. Each incident gamma and its secondaries share an event identifier, linking the detected light to the primary's interaction history. Depth selection retains all detected photons from a selected event, including light produced after later interactions elsewhere.
The entrance-side sample includes light from events whose first gamma interaction in the crystal occurred within $\qty{-5}{\milli\metre}\leq z<\qty{-4}{\milli\metre}$, while the detector-side sample uses $\qty{4}{\milli\metre}\leq z\leq\qty{5}{\milli\metre}$.


As shown in \cref{fig:depth-conditioning}, the entrance-side map is qualitatively similar to the full \qty{100}{\kilo\electronvolt} map: both are broad and smooth. Their FWHM values differ by approximately a factor of 1.5. This difference is expected because the full low-energy sample follows an exponential interaction-depth distribution, whereas the conditioned sample selects events whose first interaction occurs in the first millimetre. The similarity nevertheless supports interaction depth as an important cause of the low-energy shape.

The complete energy scan uses \qtylist{50;75;100;150;200;300;350;400;500;662;800;900;1000}{\kilo\electronvolt} to resolve the transition in more detail. The FWHM is approximately \qty{14}{\milli\metre} at \qty{50}{\kilo\electronvolt} and \qty{75}{\kilo\electronvolt}, falls to \qty{8.9}{\milli\metre} at \qty{100}{\kilo\electronvolt}, and reaches its global minimum of \qty{0.14}{\milli\metre} at \qty{200}{\kilo\electronvolt}. It then rises to approximately \qty{0.2}{\milli\metre} above \qty{662}{\kilo\electronvolt}. Interaction depth alone cannot explain this later increase. Gamma scattering provides a possible explanation and is examined in \cref{sec:interactions}. Although the FWHM changes by two orders of magnitude, the tail probability $P(r>\qty{5}{\milli\metre})$ decreases only from approximately 0.83 at low energy to 0.75 at high energy. The response therefore changes from a broad low-energy distribution into a narrow high-energy core superposed on an outer component that remains dominant.

\subsection{Influence of source geometry}
\subsubsection{Incidence-angle dependence}

In addition to gamma energy, the incidence angle can affect the optical distribution. To examine this dependence, a point beam irradiates the \ce{CeBr3} crystal with \num{20000} incident \qty{662}{\kilo\electronvolt} gammas at \ang{0}, \ang{30}, and \ang{60} from the centre of the entrance face.

The straight-line exit position is
\begin{equation}
  x_{\mathrm{exit}}=\qty{10}{\milli\metre}\tan\theta,
\end{equation}
and therefore equals approximately \qty{0}{\milli\metre}, \qty{5.8}{\milli\metre}, and \qty{17}{\milli\metre} for the three angles.

\begin{figure}[H]
  \centering
  \thesisfigure{\textwidth}{chapter6_1}{figure_6_1_angle_dependence}
  \caption{Unit-normalized detector-plane optical-photon densities for \ce{CeBr3} irradiated by a \qty{662}{\kilo\electronvolt} point beam at \ang{0}, \ang{30}, and \ang{60} (top), with the corresponding $x$ marginals integrated over all $y$ (bottom). The maps use exact \qtyproduct{1 x 1}{\milli\metre} cells and a shared logarithmic colour scale. Marginals use uniform \qty{0.25}{\milli\metre} bins from \qty{-25.5}{\milli\metre} to \qty{25.5}{\milli\metre}. Dashed lines and triangles mark $x_{\mathrm{exit}}=\qty{10}{\milli\metre}\tan\theta$.}
  \label{fig:angle-dependence}
\end{figure}

As \cref{fig:angle-dependence} shows, increasing the angle makes the distribution more asymmetric and shifts its peak. The loss of symmetry at energies of \qty{100}{\kilo\electronvolt} and above was reported qualitatively by \textcite[p. 36]{ovsiannikova2020light}.

After integration over $y$, the peaks lie at $x=\qty{-0.13(13)}{\milli\metre}$, $x=\qty{5.63(13)}{\milli\metre}$, and $x=\qty{16.88(13)}{\milli\metre}$. Peak positions are the centres of the highest \qty{0.25}{\milli\metre} marginal bins, with half-bin-width uncertainties. The inclined peaks differ from the straight-line exit by only \qty{-0.15(13)}{\milli\metre} at \ang{30} and \qty{-0.45(13)}{\milli\metre} at \ang{60}, showing that the peak closely follows the projected exit position.


\subsubsection{Finite beam width}

Real beams have a finite transverse extent. The point source is therefore compared with a newly implemented uniform disk of \qty{0.5}{\milli\metre} diameter and a circular Gaussian source of \qty{0.5}{\milli\metre} FWHM. These profiles change transverse source positions while retaining parallel particle directions. Each source irradiates \ce{CeBr3} at \qty{662}{\kilo\electronvolt} and perpendicular incidence with \num{20000} primaries.

\begin{figure}[H]
  \centering
  \thesisfigure{\textwidth}{chapter6_1}{figure_6_1_source_profiles}
  \caption{Unit-normalized radial optical-photon densities for \ce{CeBr3} irradiated at \qty{662}{\kilo\electronvolt} and \ang{0} by a point beam, a uniform disk of \qty{0.5}{\milli\metre} diameter, and a circular Gaussian beam of \qty{0.5}{\milli\metre} FWHM. The central region $\qty{0}{\milli\metre}\leq r\leq\qty{1}{\milli\metre}$ uses a linear scale and \qty{0.025}{\milli\metre} bins (left). The outer region $\qty{1}{\milli\metre}\leq r\leq\qty{25.5}{\milli\metre}$ uses a logarithmic scale and the standard outer radial bins: \qty{0.25}{\milli\metre} from \qty{1}{\milli\metre} to \qty{5}{\milli\metre}, \qty{1}{\milli\metre} from \qty{5}{\milli\metre} to \qty{25}{\milli\metre}, and a final \qty{0.5}{\milli\metre} bin. The line at $r=\qty{0.25}{\milli\metre}$ marks both the disk radius and the radial half-maximum position of the Gaussian source. Lines connect the exact bin-centre densities.}
  \label{fig:source-profiles}
\end{figure}

As shown in \cref{fig:source-profiles}, the radial FWHM values are \qty{0.23}{\milli\metre}, \qty{0.61}{\milli\metre}, and \qty{0.87}{\milli\metre} for the point, disk, and Gaussian sources. The corresponding $P(r>\qty{5}{\milli\metre})$ values are approximately \num{0.757} for all sources. A finite source therefore broadens the central peak without noticeably changing the outer component.

\subsection{Effect of gamma scattering}
\label{sec:scattering-control}

To quantify the difference between two unit-normalized maps with pixel probabilities $d_i$ and $m_i$, total variation is defined as
\begin{equation}
  \operatorname{TV}(d,m)=\frac12\sum_i|d_i-m_i|.
  \label{eq:total-variation}
\end{equation}
It is zero for identical maps.


To investigate the influence of scattering on the scintillation photon distribution, a simulation with disabled scattering is performed using \num{20000} incident \qty{662}{\kilo\electronvolt} gammas at normal incidence in the same \ce{CeBr3} geometry and optical configuration. In this scattering-disabled comparison run, gamma Compton and Rayleigh scattering are switched off. Photoelectric absorption, secondary-particle transport, scintillation, and optical transport remain enabled and unchanged.

\begin{figure}[H]
  \centering
  \thesisfigure{\textwidth}{chapter6_1}{figure_6_1_scattering_control}
  \caption{Unit-normalized radial optical-photon densities for matched \qty{662}{\kilo\electronvolt} \ce{CeBr3} simulations with gamma scattering enabled and disabled. The central region uses a linear scale with common bin widths of \qty{0.01}{\milli\metre} from \qty{0}{\milli\metre} to \qty{0.10}{\milli\metre}, \qty{0.025}{\milli\metre} from \qty{0.10}{\milli\metre} to \qty{0.25}{\milli\metre}, and \qty{0.05}{\milli\metre} from \qty{0.25}{\milli\metre} to \qty{1}{\milli\metre} (left). The outer region uses the standard radial bins on a logarithmic scale: \qty{0.25}{\milli\metre} from \qty{1}{\milli\metre} to \qty{5}{\milli\metre}, \qty{1}{\milli\metre} from \qty{5}{\milli\metre} to \qty{25}{\milli\metre}, and a final \qty{0.5}{\milli\metre} bin. Counts are divided by annular area and by the total detected-photon count. The underlying normalized \numproduct{51 x 51} maps differ by TV 0.034.}
  \label{fig:scattering-control}
\end{figure}

As \cref{fig:scattering-control} shows, enabling gamma scattering raises the mean radius from \qty{11.7}{\milli\metre} to \qty{12.3}{\milli\metre}, the median radius from \qty{10.0}{\milli\metre} to \qty{10.8}{\milli\metre}, and $P(r>\qty{5}{\milli\metre})$ from 0.72 to 0.76. The change is nevertheless not a uniform broadening. In the central bins, the scattering-enabled density is higher from the origin to about $r=\qty{0.10}{\milli\metre}$, whereas the scattering-disabled density is higher over the rest of the central panel. Gamma scattering therefore shifts normalized probability towards both the innermost core and the outer component at the expense of intermediate radii.
\FloatBarrier

% chktex-file 3
% chktex-file 25
% chktex-file 29
\section{Gamma interactions and energy deposition}
\label{sec:interactions}

The comparison with scattering disabled shows that gamma scattering broadens the outer distribution while also enhancing the innermost core. To understand these two effects, the following analysis separates displacement of the primary gamma from local energy spreading by its secondary particles.

Unless stated otherwise, this section and \cref{sec:bgo} use perpendicular irradiation at \qty{100}{\kilo\electronvolt}, \qty{662}{\kilo\electronvolt}, and \qty{1000}{\kilo\electronvolt}. Gamma-history and energy-deposition samples contain \num{100000} incident gammas per material and energy. Optical comparisons use separate \num{20000}-primary samples.


\subsection{Interaction probabilities and event classes}

Each primary is assigned to one of four mutually exclusive classes according to its first interaction in the crystal: no interaction, photoelectric first, Compton first, or Rayleigh first. This classification determines how energy is transferred and which later histories can occur.

As shown in \cref{tab:cebr-first-interactions}, Rayleigh-first histories account for approximately \qty{6}{\percent}, \qty{3}{\percent}, and \qty{2}{\percent} of interacting \ce{CeBr3} events at \qty{100}{\kilo\electronvolt}, \qty{662}{\kilo\electronvolt}, and \qty{1000}{\kilo\electronvolt}. These histories are included in the overall first-interaction statistics. However, because they are relatively infrequent at all three energies and a Rayleigh collision transfers no energy, they are discussed in less detail than Compton interactions, which dominate at high energies. At \qty{100}{\kilo\electronvolt}, gamma histories overwhelmingly consist of direct photoelectric absorption.

\begin{table}[htbp]
\centering\small
\caption{First-interaction statistics for \ce{CeBr3}. Each material--energy case contains \num{100000} incident gammas at normal incidence. The interaction probability refers to all incident gammas, while process shares are conditioned on an interaction in the scintillator.}
\label{tab:cebr-first-interactions}
\begin{tabular}{rrrrr}
\toprule
$E$ (\unit{\kilo\electronvolt}) & Interaction & Photoelectric & Compton & Rayleigh \\
\midrule
100 & 99.9\% & 85.0\% & 8.8\% & 6.2\% \\
662 & 32.6\% & 9.5\% & 87.5\% & 3.0\% \\
1000 & 26.4\% & 5.3\% & 93.0\% & 1.7\% \\
\bottomrule
\end{tabular}
\end{table}

\subsection{Repeated scattering, subsequent absorption, and escape}
\label{sec:cebr3-fates}

Because the first Compton collision is insufficient to determine the final energy deposition, Compton-first histories are classified by eventual photoabsorption or escape and by subsequent scattering counts. \Cref{tab:cebr-multiplicity} shows that eventual absorption occurs for \qty{91}{\percent}, \qty{30}{\percent}, and \qty{23}{\percent} of Compton-first \ce{CeBr3} histories at \qty{100}{\kilo\electronvolt}, \qty{662}{\kilo\electronvolt}, and \qty{1000}{\kilo\electronvolt}. The share that escapes without any additional in-crystal gamma scattering rises from \qty{8.6}{\percent} to \qty{54}{\percent} and \qty{60}{\percent}. These immediate-escape histories deposit energy only through the first Compton recoil and its local secondaries.

\begin{table}[htbp]
\centering\small
\caption{Eventual fate and subsequent scattering of Compton-first \ce{CeBr3} histories. Samples contain \num{100000} incident gammas per energy at normal incidence. Percentages refer to the corresponding Compton-first total. ``Additional Compton'' also includes histories that contain both later Rayleigh and Compton scattering.}
\label{tab:cebr-multiplicity}
\begin{tabular}{rlrrrr}
\toprule
$E$ (\unit{\kilo\electronvolt}) & Fate & Total & \shortstack{No additional\\scattering} & \shortstack{Additional\\Rayleigh only} & \shortstack{Additional\\Compton} \\
\midrule
100 & Absorbed & 90.8\% & 81.4\% & 4.1\% & 5.4\% \\
100 & Escaped & 9.2\% & 8.6\% & 0.4\% & 0.2\% \\
662 & Absorbed & 29.9\% & 15.3\% & 0.7\% & 14.0\% \\
662 & Escaped & 70.1\% & 54.3\% & 1.5\% & 14.2\% \\
1000 & Absorbed & 23.1\% & 10.7\% & 0.4\% & 12.1\% \\
1000 & Escaped & 76.9\% & 60.1\% & 1.2\% & 15.5\% \\
\bottomrule
\end{tabular}
\end{table}

Interaction frequency alone does not determine the scintillation source, because the histories can deposit different amounts of the incident gamma energy. For each history class, the fraction of incident energy deposited is defined as its summed $E_{\mathrm{dep}}$ divided by the total energy of all incident gammas. \Cref{fig:event-energy-weight} compares this fraction with the corresponding fraction of incident events. Escaped Compton-first histories are separated according to whether the primary gamma escapes without later in-crystal scattering or undergoes further Compton or Rayleigh scattering first. This distinction shows how much energy is deposited by a local first recoil and how much is associated with additional gamma transport.

\begin{figure}[!htbp]
  \captionsetup{font=small}
  \centering
  \thesisfigure{\textwidth}{chapter6_2}{figure_6_2_event_energy_weight}
  \caption{Fraction of incident events and fraction of total incident gamma energy deposited for photoelectric-first, eventually absorbed Compton-first, immediate-escape Compton-first, and Compton-first histories that escape after further primary-gamma Compton or Rayleigh interactions in \ce{CeBr3} at \qty{100}{\kilo\electronvolt}, \qty{662}{\kilo\electronvolt}, and \qty{1000}{\kilo\electronvolt}, using \num{100000} incident gammas per energy at normal incidence. Blue bars show event fractions and orange bars deposited-energy fractions. Immediate escape denotes escape without any later recorded in-crystal primary-gamma scattering. Rayleigh-first and noninteracting histories are omitted, but remain in the incident-total denominators.}
  \label{fig:event-energy-weight}
\end{figure}

At \qty{100}{\kilo\electronvolt}, the event and deposited-energy fractions are nearly equal because photoelectric-first and eventually absorbed Compton-first histories deposit almost the full incident energy. At \qty{662}{\kilo\electronvolt} and \qty{1000}{\kilo\electronvolt}, escaped Compton-first histories are much more prominent in the event count than in the deposited energy. Eventually absorbed histories do not deposit the full incident energy because secondary particles escape. Histories that terminate in the crystal contribute disproportionately to the scintillation source relative to their event frequency, while the more numerous escape histories generally leave only partial energy deposits.

\subsection{Initial Compton-scattering kinematics}

The scattering kinematics help explain the influence of Compton-first histories by linking the initial energy transfer to the direction and remaining energy of the scattered gamma. In the representative \qty{662}{\kilo\electronvolt} \ce{CeBr3} sample, the median first-scatter angle is $\chi_{50}\approx\ang{60}$. If such a collision occurred at the entrance, straight propagation across the \qty{10}{\milli\metre} thickness would correspond to a transverse displacement of $\qty{10}{\milli\metre}\tan\chi_{50}\approx\qty{17}{\milli\metre}$.

The free-electron relation in \cref{eq:compton-energy} links the scattering angle to the outgoing gamma energy and recoil-energy transfer.
The Compton model includes bound-electron effects, so this expression is only a reference~\parencite{geant4polarizedcomptonsource}. Nevertheless, at \qty{662}{\kilo\electronvolt}, the simulated angular distribution closely follows the free-electron Klein--Nishina prediction, as shown in \cref{fig:compton-kinematics}.
\Cref{fig:compton-kinematics} also shows how the initial scattering can strongly affect the eventual fate of the gamma. Gammas that scatter at $\ang{90}$ have the longest distance to travel before escaping the crystal, making them more likely to be absorbed or scatter again. At the same time, gammas that lose a large fraction of their energy during backscattering are also less likely to escape because of the higher attenuation at lower energies.

\begin{figure}[!htbp]
  \captionsetup{font=small}
  \centering
  \thesisfigure{\textwidth}{chapter6_2}{figure_6_2_compton_kinematics}
  \caption{First-Compton-angle density for \qty{662}{\kilo\electronvolt} \ce{CeBr3} in \ang{4} bins from \ang{0} to \ang{180}. The filled bars stack eventual in-crystal photoabsorption (orange), escape without later primary-gamma scattering (blue), and escape after one or more later Compton or Rayleigh collisions (green). The dashed curve is the normalized free-electron Klein--Nishina expectation from \cref{eq:klein-nishina-angle}. The upper axis gives the free-electron outgoing energy $E'(\chi)$ at the angles marked on the lower axis.}
  \label{fig:compton-kinematics}
\end{figure}

\subsection{Transverse gamma displacement}

Energy can be deposited away from the beam because the primary gamma scatters and interacts again at another position, or because a charged secondary particle or fluorescence photon travels away from the interaction point. To separate these effects, the transverse displacement of the primary gamma is measured from its first interaction to its eventual absorption point. For Compton-first and Rayleigh-first histories that end in photoelectric absorption, this displacement is
\begin{equation}
  \Delta r_\perp=\sqrt{(x_{\mathrm{abs}}-x_{\mathrm{first}})^2+(y_{\mathrm{abs}}-y_{\mathrm{first}})^2}.
\end{equation}
\Cref{fig:gamma-displacement} shows the Compton-first distribution of this displacement. The Compton-first medians in \ce{CeBr3} are approximately \qty{0.5}{\milli\metre} at \qty{100}{\kilo\electronvolt}, \qty{4.6}{\milli\metre} at \qty{662}{\kilo\electronvolt}, and \qty{5.4}{\milli\metre} at \qty{1000}{\kilo\electronvolt}. These lengths are substantial relative to the \qty{25.5}{\milli\metre} half-width. This displacement provides a mechanism for the scattering-induced shift towards the outer component observed in \cref{sec:scattering-control}.

\begin{figure}[!htbp]
  \captionsetup{font=small}
  \centering
  \thesisfigure{\textwidth}{chapter6_2}{figure_6_2_gamma_displacement}
  \caption{Transverse displacement from the first interaction to eventual primary-gamma photoabsorption for absorbed Compton-first \ce{CeBr3} histories, drawn from \num{100000} incident gammas per energy at perpendicular incidence, at \qty{100}{\kilo\electronvolt}, \qty{662}{\kilo\electronvolt}, and \qty{1000}{\kilo\electronvolt}. The \qty{100}{\kilo\electronvolt} panel uses \qty{0.1}{\milli\metre} bins over \qtyrange{0}{3}{\milli\metre}, while the higher-energy panels use \qty{0.5}{\milli\metre} bins below \qty{5}{\milli\metre} and \qty{1}{\milli\metre} bins up to \qty{25}{\milli\metre}. Counts are divided by the full absorbed Compton-first count and bin width, without renormalizing to the displayed range. Dotted and dashed vertical lines mark the median and 90th percentile of the full displacement sample.}
  \label{fig:gamma-displacement}
\end{figure}


\subsection{Local spreading by secondary particles}

Aside from gamma transport, secondary particles can also carry energy away from an interaction point. Photoelectric-first events provide one measure of local secondary-particle spreading: the primary gamma terminates at the first interaction, while secondary particles carry energy away from that point. A Compton interaction may behave differently from a photoelectric interaction in this regard due to the lower energy of an ejected Compton-recoil electron. To examine this contribution, Compton-first histories that escape without any further in-crystal interaction are considered. Their deposited energy comes from the first recoil and its secondaries, with no contribution from a later gamma interaction at a displaced position.

Local spreading is quantified by the radius enclosing half of a history class's summed deposited energy, measured transversely from the first interaction point. For photoelectric-first events, these radii are \qty{0.0043}{\milli\metre}, \qty{0.14}{\milli\metre}, and \qty{0.23}{\milli\metre} at \qty{100}{\kilo\electronvolt}, \qty{662}{\kilo\electronvolt}, and \qty{1000}{\kilo\electronvolt}. For immediate-escape Compton-first histories, the corresponding medians are \qty{0.00012}{\milli\metre}, \qty{0.039}{\milli\metre}, and \qty{0.091}{\milli\metre}. The local Compton-recoil branch is therefore more concentrated than the photoelectric branch at all three energies, providing a mechanism for the enhanced innermost core observed in \cref{sec:scattering-control}.

\subsection{Explanation of the moderate broadening caused by gamma scattering}

Gamma scattering produces two spatially different contributions. At its first Compton collision, a gamma creates a recoil electron whose deposited energy is more localized than in photoelectric-first events. By contrast, gammas that interact again after their initial scatter move deposited energy several millimetres away from the beam axis and feed the outer component.

Gamma scattering therefore changes both the core and tail in \cref{fig:scattering-control}. Because optical transport dominates the overall distribution, however, gamma scattering has only a moderate effect.
\FloatBarrier

% chktex-file 29
\section{BGO}
\label{sec:bgo}

\subsection{Motivation, implementation, and expected physical consequences}

A material with a larger photoelectric contribution is expected to reduce the spatial redistribution caused by scattered primary gammas. For this reason, high-$Z$, high-density \ce{Bi4Ge3O12} (BGO) is of interest and was added as a selectable material to the simulation. It is compared with \ce{CeBr3} using the same crystal dimensions, reflector, housing, beam, and detection plane. The relevant configured properties are summarized in \cref{tab:material-properties}. BGO's bismuth content and density are expected to provide stronger attenuation and a larger photoelectric contribution.

This combination is expected to produce a more compact deposited-energy distribution by reducing the average transverse displacement of primary gammas.

\begin{table}[htbp]
\centering\small
\caption{Scintillator properties configured in the simulation. Yields give the mean number of emitted photons per unit deposited energy. Bulk optical absorption is the optical-photon absorption length, with the \ce{CeBr3} range reflecting its dependence on optical-photon energy. The BGO model adopts the \qty{4.00}{\metre} absorption-length input used in the \qty{480}{\nano\metre} optical simulations of \textcite[section II-B and table I]{valenciaga2018geometry}, applying it here independently of wavelength.}
\label{tab:material-properties}
\begin{tabularx}{\textwidth}{l>{\raggedright\arraybackslash}X>{\raggedright\arraybackslash}X>{\raggedright\arraybackslash}X}
\toprule
Material & Density & \shortstack{Scintillation\\yield} & \shortstack{Optical absorption\\length} \\
\midrule
\ce{CeBr3} & \qty{5.20}{\gram\per\cubic\centi\metre} & \qty{60}{\per\kilo\electronvolt}\newline\parencite{ovsiannikova2020light} & \qtyrange{105}{155}{\milli\metre}, energy dependent \\
\ce{Bi4Ge3O12} & \qty{7.13}{\gram\per\cubic\centi\metre}\newline\parencite{miller2025material} & \qty{8.5}{\per\kilo\electronvolt}\newline\parencite{miller2025material} & \qty{4.00}{\metre}, constant\newline\parencite{valenciaga2018geometry} \\
\bottomrule
\end{tabularx}
\end{table}


\subsection{Interaction probabilities and event classes compared with \ce{CeBr3}}

The material comparison uses \qty{100}{\kilo\electronvolt}, \qty{662}{\kilo\electronvolt}, and \qty{1000}{\kilo\electronvolt}, while all other settings follow the defaults defined in \cref{chap:simulations}.
At \qty{662}{\kilo\electronvolt}, approximately \qty{51}{\percent} of incident gammas interact in BGO, compared with \qty{33}{\percent} in \ce{CeBr3}. At that energy, among interacting primaries, the BGO sample is \qty{32}{\percent} photoelectric-first and \qty{64}{\percent} Compton-first, compared with \qty{9.5}{\percent} and \qty{88}{\percent} in \ce{CeBr3}. The three-energy summary is given in \cref{tab:material-first-interactions}. This confirms the expected increase in overall interaction probability and in the fraction of photoelectric interactions in BGO.

\begin{table}[htbp]
\centering\small
\caption{First-interaction statistics for \ce{CeBr3} and BGO. Each material--energy case contains \num{100000} incident gammas at normal incidence. The interaction probability refers to all incident gammas, while process shares are conditioned on an interaction in the scintillator. Independently rounded process shares may not sum to exactly \qty{100}{\percent}.}
\label{tab:material-first-interactions}
\begin{tabular}{lrrrrr}
\toprule
Material & $E$ (\unit{\kilo\electronvolt}) & Interaction & Photoelectric & Compton & Rayleigh \\
\midrule
\ce{CeBr3} & 100 & 99.9\% & 85.0\% & 8.8\% & 6.2\% \\
BGO & 100 & 100.0\% & 93.8\% & 2.9\% & 3.4\% \\
\ce{CeBr3} & 662 & 32.6\% & 9.5\% & 87.5\% & 3.0\% \\
BGO & 662 & 51.0\% & 31.5\% & 64.1\% & 4.4\% \\
\ce{CeBr3} & 1000 & 26.4\% & 5.3\% & 93.0\% & 1.7\% \\
BGO & 1000 & 38.6\% & 19.4\% & 77.7\% & 3.0\% \\
\bottomrule
\end{tabular}
\end{table}

\subsection{Fate and displacement of scattered gammas}

Following the classification in \cref{sec:cebr3-fates}, Compton-first histories are separated by eventual photoabsorption or escape and by subsequent scattering. \Cref{tab:material-compton-fates} gives the BGO fractions. The eventually absorbed fractions of Compton-first primaries in BGO are \qty{79}{\percent}, \qty{59}{\percent}, and \qty{48}{\percent} at \qty{100}{\kilo\electronvolt}, \qty{662}{\kilo\electronvolt}, and \qty{1000}{\kilo\electronvolt}, compared with \qty{91}{\percent}, \qty{30}{\percent}, and \qty{23}{\percent} in \ce{CeBr3}. The higher absorbed fraction in BGO at high energies is expected, but the reversal at \qty{100}{\kilo\electronvolt} requires explanation.

The lower eventual absorption in BGO cannot be attributed to a difference in Compton kinematics, which are nearly independent of the material. \Cref{fig:material-compton-angles} illustrates this weak material dependence. At \qty{662}{\kilo\electronvolt}, the mean first-scatter angles are nearly identical: \ang{68.3} in \ce{CeBr3} and \ang{68.0} in BGO. A shallower first-Compton interaction depth provides an explanation for the lower absorbed fraction at \qty{100}{\kilo\electronvolt}. At this energy, the median depth of the first-Compton interaction in BGO is \qty{0.25}{\milli\metre} compared with \qty{1.03}{\milli\metre} in \ce{CeBr3}. Consequently, after a gamma backscatters towards the entrance, it has a shorter path to escape in BGO than in \ce{CeBr3}.

\begin{figure}[!htbp]
  \captionsetup{font=small}
  \centering
  \thesisfigure{\textwidth}{chapter6_3}{figure_6_3_material_compton_angles}
  \caption{Unit-normalized first-Compton-angle densities for \ce{CeBr3} and BGO at \qty{662}{\kilo\electronvolt}, using common \ang{4} bins from \ang{0} to \ang{180}. Filled orange bars show BGO, and an unfilled blue step histogram shows \ce{CeBr3} with identical bin boundaries. The dashed curve is the normalized free-electron Klein--Nishina expectation from \cref{eq:klein-nishina-angle}. The upper axis gives the free-electron outgoing energy $E'(\chi)$.}
  \label{fig:material-compton-angles}
\end{figure}

\begin{table}[htbp]
\centering\small
\caption{Eventual fate and subsequent scattering of Compton-first BGO histories. Samples contain \num{100000} incident gammas per energy at normal incidence. Percentages refer to the corresponding Compton-first total. ``Additional Compton'' includes histories that also contain later Rayleigh scattering. Independently rounded entries may not add to the displayed totals.}
\label{tab:material-compton-fates}
\begin{tabular}{rlrrrr}
\toprule
$E$ (\unit{\kilo\electronvolt}) & Fate & Total & \shortstack{No additional\\scattering} & \shortstack{Additional\\Rayleigh only} & \shortstack{Additional\\Compton} \\
\midrule
100 & Absorbed & 78.7\% & 70.7\% & 4.7\% & 3.2\% \\
100 & Escaped & 21.3\% & 19.5\% & 1.0\% & 0.8\% \\
662 & Absorbed & 59.2\% & 39.8\% & 1.7\% & 17.6\% \\
662 & Escaped & 40.8\% & 33.9\% & 1.3\% & 5.6\% \\
1000 & Absorbed & 48.4\% & 30.0\% & 1.3\% & 17.1\% \\
1000 & Escaped & 51.6\% & 43.0\% & 1.0\% & 7.6\% \\
\bottomrule
\end{tabular}
\end{table}

\Cref{fig:material-displacements} compares the transverse displacement distribution of absorbed Compton-first histories in BGO and \ce{CeBr3}. It confirms the expectation that scattered gammas travel a shorter distance in BGO, with median transverse displacements in \ce{CeBr3}/BGO of \qty{0.48}{\milli\metre}/\qty{0.23}{\milli\metre} at \qty{100}{\kilo\electronvolt}, \qty{4.6}{\milli\metre}/\qty{2.1}{\milli\metre} at \qty{662}{\kilo\electronvolt}, and \qty{5.4}{\milli\metre}/\qty{2.7}{\milli\metre} at \qty{1000}{\kilo\electronvolt}. The medians are therefore about a factor of two smaller in BGO.



\begin{figure}[!htbp]
  \captionsetup{font=small}
  \centering
  \thesisfigure{\textwidth}{chapter6_3}{figure_6_3_material_displacements}
  \caption{Transverse primary-gamma displacement from first interaction to eventual photoabsorption for absorbed Compton-first \ce{CeBr3} and BGO histories at \qty{100}{\kilo\electronvolt}, \qty{662}{\kilo\electronvolt}, and \qty{1000}{\kilo\electronvolt}. Each material--energy case contains \num{100000} incident gammas at normal incidence. Each curve is normalized by the full absorbed Compton-first count and bin width, not just the displayed range. Blue denotes \ce{CeBr3} and orange BGO. Each energy has its own linear vertical scale. The \qty{100}{\kilo\electronvolt} panel uses \qty{0.1}{\milli\metre} bins over \qtyrange{0}{3}{\milli\metre}, while the other panels use \qty{0.5}{\milli\metre} bins below \qty{5}{\milli\metre} and \qty{1}{\milli\metre} bins up to \qty{25}{\milli\metre}.}
  \label{fig:material-displacements}
\end{figure}

\subsection{Energy-deposition comparison with \ce{CeBr3}}
\label{sec:material-edep-method}

The material-dependent gamma histories translate directly into the deposited-energy field shown in \cref{fig:material-edep}. Both materials are compared on the same radial and depth grids, with narrow bins near the beam axis and entrance face and wider bins in the sparsely populated outer regions.\footnote{\label{fn:material-edep-bins}Radial boundaries in millimetres are 0, 0.0025, 0.005, 0.01, 0.02, 0.05, 0.1, 0.2, 0.3, 0.5, 0.75, 1, 1.5, 2, 3, 5, 7.5, 10, 15, 20, and 25.5. Depth bins have widths of \qty{0.1}{\milli\metre} from 0 to \qty{1}{\milli\metre}, \qty{0.2}{\milli\metre} from 1 to \qty{2}{\milli\metre}, then boundaries at 2.5, 3, 4, 6, and \qty{10}{\milli\metre}.}

At \qty{662}{\kilo\electronvolt}, BGO deposits \qty{258}{\kilo\electronvolt} per incident primary, roughly twice the \qty{123}{\kilo\electronvolt} in \ce{CeBr3}. Only \qty{6.2}{\percent} of this energy lies beyond \qty{5}{\milli\metre}, compared with \qty{11}{\percent} in \ce{CeBr3}. The same ordering is found at \qty{1000}{\kilo\electronvolt}. BGO's larger photoelectric-first share, stronger reabsorption of scattered primaries, and shorter transverse gamma displacement therefore reduce the outer part of the deposited-energy distribution.

\begin{figure}[!htbp]
  \captionsetup{font=small}
  \centering
  \thesisfigure{0.94\textwidth}{chapter6_3}{figure_6_3_material_edep}
  \caption{Normalized deposited-energy densities for \ce{CeBr3} (blue) and BGO (orange), using \num{100000} incident gammas per material and energy at perpendicular incidence. Bin energies are divided by the full in-crystal deposited energy and by radial or depth bin width, respectively. The top panels compare radii from the beam axis at \qty{100}{\kilo\electronvolt}, \qty{662}{\kilo\electronvolt}, and \qty{1000}{\kilo\electronvolt}, with a common logarithmic vertical scale and a linear radial axis spanning \qtyrange{0}{25.5}{\milli\metre}. All radial curves use one common variable-width radial grid, which is narrow near the beam axis and broader in the sparse outer region. The bottom panel overlays entrance-referenced depth densities at \qty{100}{\kilo\electronvolt} (solid) and \qty{662}{\kilo\electronvolt} (dashed), using material colours and a logarithmic vertical scale. All depth curves use a separate variable-width depth grid, which resolves the entrance region more finely than the detector-side region. Both radial and depth bin boundaries are specified in \hyperref[fn:material-edep-bins]{footnote~\ref*{fn:material-edep-bins}}.}
  \label{fig:material-edep}
\end{figure}

At \qty{100}{\kilo\electronvolt}, the width ordering reverses. The beam-axis radius $r_{50}$ encloses half of the total deposited energy. BGO has $r_{50}=\qty{0.039}{\milli\metre}$, compared with $r_{50}=\qty{0.0053}{\milli\metre}$ in \ce{CeBr3}, and deposits \qty{86}{\kilo\electronvolt} rather than \qty{99}{\kilo\electronvolt} per primary.
Nearly all incident gammas interact in both materials at this energy. In BGO, however, interactions occur closer to the entrance and more of the resulting energy escapes from that nearby surface. This removes part of the strongly beam-centred deposition. After the remaining deposited energy is normalized to one, energy deposited farther from the axis consequently has a larger relative weight.

\begin{figure}[!t]
  \captionsetup{font=small}
  \centering
  \thesisfigure{\textwidth}{chapter6_3}{figure_6_3_material_optical_shapes}
  \caption{Unit-normalized \ce{CeBr3} and BGO detector maps under perpendicular point-beam irradiation, with \num{20000} incident gammas per material and energy, at \qty{100}{\kilo\electronvolt}, \qty{662}{\kilo\electronvolt}, and \qty{1000}{\kilo\electronvolt}, using uniform \qtyproduct{1 x 1}{\milli\metre} cells and one common logarithmic probability-density scale (top two rows). The bottom row displays the corresponding radial probability densities using the standard optical radial bins specified in \cref{fig:energy-maps}.}
  \label{fig:material-optical-shapes}
\end{figure}

\subsection{Normalized optical-photon distribution compared with \ce{CeBr3}}

The optical-photon distributions in \cref{fig:material-optical-shapes}, obtained using \num{20000} primaries per material and energy, show the same qualitative energy dependence in both materials. At \qty{100}{\kilo\electronvolt}, \qty{662}{\kilo\electronvolt}, and \qty{1000}{\kilo\electronvolt}, the radial FWHM values are \qty{8.9}{\milli\metre}, \qty{0.23}{\milli\metre}, and \qty{0.22}{\milli\metre} in \ce{CeBr3}, and \qty{13.8}{\milli\metre}, \qty{0.22}{\milli\metre}, and \qty{0.24}{\milli\metre} in BGO. The corresponding $P(r>\qty{5}{\milli\metre})$ values are 0.83, 0.76, and 0.75 for \ce{CeBr3}, compared with 0.85, 0.78, and 0.77 for BGO. BGO therefore has a larger tail probability at all three energies and a larger FWHM at \qty{100}{\kilo\electronvolt} and \qty{1000}{\kilo\electronvolt}, while at \qty{662}{\kilo\electronvolt} its FWHM is marginally smaller.

BGO's longer bulk absorption length allows more optical photons following long or reflected paths to survive. At the same time, its stronger gamma attenuation moves the scintillation source towards the entrance face, increasing the distance to the detector plane. These effects enlarge the broad detector-plane component even though the high-energy deposited-energy field is more compact.

\subsection{Absolute optical-photon yield compared with \ce{CeBr3}}

One major material difference is scintillation yield: the configured BGO-to-\ce{CeBr3} yield ratio is $8.5/60\approx 0.14$. The corresponding detected-photon count ratios are 0.15, 0.34, and 0.31 at \qty{100}{\kilo\electronvolt}, \qty{662}{\kilo\electronvolt}, and \qty{1000}{\kilo\electronvolt}. The higher BGO interaction probability and longer optical absorption length partly compensate for its smaller yield at high energies. At \qty{100}{\kilo\electronvolt}, nearly every incident gamma produces detected light in both materials and the ratio remains close to the configured yield ratio.

\subsection{Assessment of material-dependent trade-offs}

For the simulated \qtyproduct{51 x 51 x 10}{\milli\metre} geometry, BGO does not provide the hoped-for sharper normalized optical response. It also produces markedly fewer detected photons. Thus, for spatial localization in this geometry, BGO offers no consistent advantage in the shape of the optical distribution and incurs a substantial signal penalty.

\FloatBarrier

% chktex-file 3
% chktex-file 12
% chktex-file 25
% chktex-file 29
\section{Parametric fits of the optical-photon distributions}
\label{sec:fits}

\subsection{Modelling target and selection criteria}
\label{sec:shape-validation}

This section examines parametric descriptions of the detector-plane distribution, which might allow beam properties such as the gamma spectrum to be inferred from a measured spatial response. At each of the 13 energies listed in \cref{sec:cebr-energy-scan}, the fit target is the complete wavelength-integrated $51\times51$ map for a point beam, normalized to unit probability. Absolute light yield is excluded. Radial and central-strip profiles visualize these full-map fits.

Model pixel probabilities $m_i$ are obtained by integrating the model over each pixel and normalizing over the detector. Fits minimize cross-entropy $L=-\sum_i d_i\ln m_i$, where $d_i$ is the observed pixel probability.

Complete events are assigned to five equally sized groups by their identifier. Each fit uses four groups and is tested on the independently normalized fifth, preventing correlated photons from one gamma from appearing in both samples. Each case is summarized by the mean of its five held-out TV values. Medians and worst values across cases refer to these case means. A useful model should combine low held-out TV as defined in \cref{eq:total-variation} with a small number of parameters.

\subsection{Candidate models for perpendicular incidence}

\subsubsection{Gaussian}
The approximately Gaussian detector-plane distributions reported by \textcite[p.~40]{ovsiannikova2020light} for \ce{CeBr3} at gamma energies up to \qty{100}{\kilo\electronvolt} motivate the one-parameter Gaussian as a starting candidate:
\begin{equation}
  G(r;\sigma)=\exp\!\left(-\frac{r^2}{2\sigma^2}\right).
\end{equation}
The Gaussian model's median held-out TV is approximately 0.2 for both materials. As shown in \cref{fig:gaussian-moffat}, its single rapidly decaying scale cannot simultaneously capture the central concentration and the broad wings.

\subsubsection{Moffat}

To allow broader wings than the Gaussian, the next candidate is a radial Moffat profile with two parameters:
\begin{equation}
  M(r;p,w)=\left(1+\frac{r^2}{pw^2}\right)^{-p}.
\end{equation}
Here, $w$ is a width scale and $p$ is a dimensionless shape exponent.
In the limit $p\to\infty$, it approaches $\exp(-r^2/w^2)$ and therefore contains a Gaussian with $\sigma=w/\sqrt{2}$ as a limiting case. Finite values of $p$ allow heavier tails.

Its median held-out TV reaches 0.05 for \ce{CeBr3} and 0.04 for BGO, but the finite core misses the sharp high-energy peak visible in \cref{fig:gaussian-moffat}.


\begin{figure}[!htbp]
  \centering
  \captionsetup{font=small}
  \thesisfigure{\textwidth}{chapter6_4}{figure_6_4_gaussian_moffat}
  \caption{Gaussian and Moffat fits to perpendicular \ce{CeBr3} profiles at 100, 662, and \qty{1000}{\kilo\electronvolt}, with \num{20000} primaries per case. Black lines with markers every third bin show photon counts divided by the complete-map count and annular area. Coloured curves show annular averages of models fitted to the complete detector map with \qtyproduct{1 x 1}{\milli\metre} cells. Radial bin widths are \qty{0.05}{\milli\metre} below \qty{1}{\milli\metre}, \qty{0.25}{\milli\metre} from 1 to \qty{5}{\milli\metre}, and \qty{1}{\milli\metre} above this, ending with a half-width bin at \qty{25.5}{\milli\metre}. All panels share the same axes, with logarithmic density.}
  \label{fig:gaussian-moffat}
\end{figure}
\FloatBarrier

\subsubsection{Direct exponential-depth kernel}

The Moffat profile's insufficient core motivates a model in which the peak arises from scintillation sites close to the detector. This can happen in the simple case of an exponentially attenuated source along the incident beam, with direct isotropic propagation from each source point to the detector plane.

For an on-axis point at normal distance $h$ from the detector, projection of the solid angle onto the plane gives the density
\begin{equation}
  K(r\mid h)\propto\frac{h}{(r^2+h^2)^{3/2}}.
\end{equation}

For a crystal of thickness $T=\qty{10}{\milli\metre}$, let $s$ denote the depth measured from the entrance. The source density is proportional to $e^{-s/\lambda}$ for $0\leq s\leq T$, giving the direct kernel
\begin{equation}
  L(r;\lambda)=A_\lambda\int_0^T
  e^{-s/\lambda}K(r\mid T-s)\,\mathrm ds.
\end{equation}
The factor $A_\lambda$ normalizes the model on the finite detector.

The median held-out TV is 0.11 for \ce{CeBr3} and 0.12 for BGO. Thus, the direct kernel performs worse than the Moffat profile over the complete map. As shown in \cref{fig:perpendicular-models}, this approach fails to reproduce both the central peak and the outer distribution. This is likely due to the effects of reflection, gamma scattering, and secondary-particle spreading.

\subsubsection{Power law}

Since the direct kernel does not describe the complete distribution, the central region is examined separately to identify a profile that can reproduce the sharp high-energy peak. In the double-logarithmic insets of \cref{fig:perpendicular-models}, the radial density is approximately linear over the resolved sub-millimetre range. This behaviour motivates the Power profile:
\begin{equation}
  P(r;q)=\left(\frac{r}{r_*}\right)^{-q},\qquad r_*=\qty{1}{\milli\metre}.
\end{equation}
The dimensionless exponent $q$ controls the radial decay.
Because the model is integrated over finite two-dimensional pixels, the central integral is finite exactly when $q<2$. Fits use $0.01\leq q\leq1.98$. The fixed scale $r_*$ only sets the units and cancels after normalization.

Its median held-out TV over the complete map is 0.09 for \ce{CeBr3} and BGO. The ability of the power law to describe the central peak is tested with a separate fit restricted to $r<\qty{1}{\milli\metre}$. This domain contains the 20 native \qty{0.05}{\milli\metre} radial bins. The median TV over all 13 energies in this domain is 0.013 for \ce{CeBr3} and 0.011 for BGO, with worst values of 0.041 and 0.027. The restricted fits shown in \cref{fig:perpendicular-models} therefore confirm that a power law describes the normalized sub-millimetre peak shape. It cannot, however, describe the broad outer component at the same time, explaining its high TV when fitted to the complete map.

\begin{figure}[!htbp]
  \centering
  \captionsetup{font=small}
  \thesisfigure{\textwidth}{chapter6_4}{figure_6_4_perpendicular_models}
  \caption{Direct exponential-depth and Power fits to the same perpendicular \ce{CeBr3} profiles as in \cref{fig:gaussian-moffat}, using the same bins and complete-map normalization. Main-panel curves retain the parameters fitted to the complete \qtyproduct{1 x 1}{\milli\metre} detector map. The insets use logarithmic radius from 0.025 to \qty{1}{\milli\metre} and compare the fine-bin data with the direct kernel, the unrestricted Power fit, and an independent Power fit restricted to the 20 bins with $r<\qty{1}{\milli\metre}$. This restricted fit is normalized to the total probability in its domain.}
  \label{fig:perpendicular-models}
\end{figure}
\FloatBarrier

\subsubsection{Selected Power--Moffat mixture}

Because the Power and Moffat candidates describe complementary features, they are combined in the mixture below. A hat denotes a component normalized separately over the finite detector.
\begin{equation}
  F_0(r)=(1-f_M)\widehat P(r;q)+f_M\widehat M(r;p,w).
\end{equation}
The weight $0\leq f_M\leq1$ is the Moffat fraction of the normalized model probability. The four-parameter mixture has a median held-out TV of approximately 0.020 for \ce{CeBr3} and 0.018 for BGO, with worst values of 0.023 and 0.021.
This mixture has a low TV, uses relatively few parameters, and describes the central concentration and broad wings simultaneously at the \qtyproduct{1 x 1}{\milli\metre} fit resolution. It is therefore selected for perpendicular incidence. \Cref{fig:selected-perpendicular-fit} illustrates the selected mixture for both materials at \qty{662}{\kilo\electronvolt}.

\begin{figure}[H]
  \centering
  \captionsetup{font=small}
  \thesisfigure{\textwidth}{chapter6_4}{figure_6_4_selected_perpendicular_fit}
  \caption{Selected Power--Moffat fits for \ce{CeBr3} and BGO at \qty{662}{\kilo\electronvolt} and perpendicular incidence, with \num{20000} primaries per case. Black points show radial data normalized by the complete-map photon count, and blue lines show the annularly averaged selected fits. Both panels use the bins and normalization of \cref{fig:perpendicular-models}, identical radius limits from 0 to \qty{25.5}{\milli\metre}, and a shared logarithmic density scale.}
  \label{fig:selected-perpendicular-fit}
\end{figure}

\subsection{Inclined-beam extensions}

The perpendicular model provides the baseline for an extension to inclined beams.

The inclined study uses 18 matched cases: two materials, three energies (\qty{100}{\kilo\electronvolt}, \qty{662}{\kilo\electronvolt}, and \qty{1000}{\kilo\electronvolt}), and three angles (\ang{0}, \ang{30}, and \ang{60}), each with \num{20000} primaries. Twelve cases have nonzero inclination.

\subsubsection{Shared-centre model}

The first extension shifts both components to a shared centre along $x$:
\begin{equation}
  \rho_P^2=\rho_M^2=(x-x_c)^2+y^2,
\end{equation}
where $\rho$ is the distance from the shared component centre. Only $x_c$ is fitted because the beam lies in the $x$--$z$ plane and the geometry remains symmetric about $y=0$. A fitted $y_c$ coordinate could be added in the same way.
This model yields a median held-out TV of 0.066 and a worst value of 0.103 because a radial profile about one centre cannot reproduce the strong asymmetry at high inclinations.

\subsubsection{Two-centre model}

To represent the asymmetry, the Power and Moffat components are instead assigned separate centres:
\begin{equation}
  \rho_P^2=(x-x_P)^2+y^2,
  \qquad \rho_M^2=(x-x_M)^2+y^2.
\end{equation}
This allows the model to represent simultaneously the sharp exit-side peak and the broad entrance-side shoulder.
The two-centre model reduces the median held-out TV to 0.033 and the worst value to 0.060. However, \cref{fig:coordinate-extensions} reveals a structural issue with the two-centre model: it generates a local minimum between the two fitted centres, while the data show a monotonic ridge. The interaction-depth dependence examined in \cref{sec:cebr-energy-scan} helps explain this discrepancy. Interactions near the entrance create a broad distribution, represented by the Moffat component, while interactions near the exit create a narrow distribution, represented by the Power component. Interactions along the beam produce a continuous transition between these shapes, which the two-centre model does not capture.

A future extension that provides a continuous ridge between the broad shoulder and the narrow peak could reduce this artificial minimum.

\begin{figure}[!htbp]
  \centering
  \captionsetup{font=small}
  \thesisfigure{\textwidth}{chapter6_4}{figure_6_4_coordinate_extensions}
  \caption{Central-strip profiles for \ce{CeBr3} at \qty{662}{\kilo\electronvolt} and incidence angles of \ang{0} and \ang{60}, with \num{20000} primaries per case. Black steps show photons in $|y|<\qty{0.5}{\milli\metre}$, binned in \qty{0.25}{\milli\metre} along $x$ and divided by the complete-map photon count and bin width. Grey, orange, and blue steps show the origin-constrained, shared-centre, and two-centre models integrated over the same \qty{0.25}{\milli\metre} strip bins. Both panels span $-25.5\leq x\leq25.5$ mm and share a logarithmic density scale. Vertical references mark the origin and the fitted centres of the shared-centre and two-centre models.}
  \label{fig:coordinate-extensions}
\end{figure}

\chapter{Conclusions and outlook}
\section{Summary}

\subsection{Physical origin of the observed shape}

The optical photon distribution is mainly determined by the interaction depth and subsequent optical transport. Low-energy gammas deposit energy near the entrance and produce a broad response. At higher energies, deposits near the readout add a sharp central peak, while the outer component remains dominant in photon count. Gamma scattering has a comparatively small effect on the complete normalized map, despite its frequent occurrence, especially at high energies. It nevertheless enhances both the innermost core, through localized recoil deposition, and the tails, through later interactions at displaced positions.

The dominant first-interaction process changes from photoelectric absorption at low energies to Compton scattering at high energies. Escape after Compton scattering becomes more frequent with increasing energy.

BGO has a larger share of photoelectric interactions and stronger reabsorption of scattered gammas. Its shorter gamma displacements make the deposited-energy distribution more compact. This does not produce a consistent optical-shape advantage because stronger attenuation near the entrance and greater survival of long optical paths increase the broad light component.

\subsection{Parametric representation}

The four-parameter Power--Moffat mixture describes the sharp central concentration and broad wings at perpendicular incidence, with median held-out TV values of 0.020 for \ce{CeBr3} and 0.018 for BGO across thirteen energies. This agreement applies at the one-millimetre fit resolution and does not guarantee an accurate sub-millimetre peak height. For inclined beams, the two-centre model achieves a median held-out TV of 0.033, but leaves an artificial minimum between the shoulder and peak where the data show a continuous ridge.

\section{Limitations}
\subsection{Simulation model}

The continuous readout omits pixelation, wavelength-dependent detection efficiency, and electronics. The optical distribution also depends on assumed optical properties, particularly the inherited \ce{LaBr3} attenuation surrogate for \ce{CeBr3} and the wavelength-independent BGO absorption length.

Constant scintillation yields exclude nonproportionality and quenching. Such effects could change the relative light contributions of different deposits and interaction histories.

\subsection{Dataset and coverage}

The results apply to one crystal geometry and selected beam conditions. Numerical widths and fitted parameters cannot be transferred directly to other detector configurations. The maximum simulated energy of \qty{1000}{\kilo\electronvolt} also excludes pair production, so its effect on containment and the optical core-to-tail balance has not been tested.

\subsection{Interpretation and validation}

The parametric models describe optical maps at specific beam energies and angles. They have not yet been used to reconstruct unknown beam parameters. Whether these descriptions also reproduce measured detector signals remains an open experimental question.

\section{Outlook}

A direct extension of the inclined model would be a continuous connection between the shoulder and peak. The existing inclined samples provide a basis for testing such a modification without changing the detector simulation. Independently of this improvement, the current models could serve as a starting point for beam reconstruction. This would first require establishing how their shape parameters depend jointly on gamma energy and incidence angle. Tests under beam conditions not used to establish this relationship could then determine which beam parameters can be recovered from spatial shape alone and where different conditions remain indistinguishable.

Further simulations could test how the observed behaviour extends beyond the present crystal geometry and energy range. Varying the crystal thickness would change both the available interaction depths and the optical propagation distances, allowing their influence on the central peak and broad component to be studied systematically. Extending the gamma energy range above the pair-production threshold would also test how the additional charged-particle tracks and annihilation photons affect the optical distribution.

Experimental validation would be needed to establish whether the simulated dependence of the optical distribution on gamma energy and incidence angle is also observed in a physical detector. A comparison with measurements could test the predicted central peak and tails and help assess the optical-property assumptions. This would require accounting for the response of the detector used in the comparison.


\appendix
% chktex-file 29

\chapter{Running further simulations}
\label{app:running-simulations}
\lstset{basicstyle=\ttfamily\footnotesize,columns=fullflexible,
  keepspaces=true,showstringspaces=false}

This appendix describes how to perform a simulation case. The source code, installation scripts, and further setup instructions are available in the simulation repository:
\begin{center}
  \url{https://github.com/Lantaan/scintillator-sim}
\end{center}

\section{Preparing the application}

In a Bash terminal on Linux or WSL Ubuntu on Windows, obtain the code and build it with the \geantfour{} version used in this work:
\begin{lstlisting}[language=bash,numbers=none]
git clone https://github.com/Lantaan/scintillator-sim.git
cd scintillator-sim
bash scripts/setup_wsl_geant4.sh \
  --geant-version 11.4.1 --without-qt --install-deps
\end{lstlisting}
The script installs \geantfour{}, its datasets, and the application. Omit \texttt{--install-deps} if system dependencies are already installed. Otherwise this step requires admin privileges. More detailed setup and troubleshooting instructions are in \path{docs/}.

For an existing installation, enter the simulation repository and rebuild after source changes with:
\begin{lstlisting}[language=bash,numbers=none]
bash scripts/build_app.sh
\end{lstlisting}
Both the build and batch scripts load the \geantfour{} environment from \path{$HOME/geant-install/bin/geant4.sh}.

\section{Running a simulation with a macro}

Simulation cases are defined as \path{.mac} files, containing the text commands to configure the application. Most importantly it specifies the output records while also containing declarations of the detector geometry, material, and incident beam. The macro can be edited without recompiling the application.

The example macro \path{default.mac} in the repository's root directory contains comments explaining the exact configuration options and can be used as a reference when preparing or modifying a macro.

From a Bash terminal in the \texttt{scintillator-sim} repository, run a specified macro and write its output into an existing directory with:

\begin{lstlisting}[language=bash,numbers=none]
OPNOVICE_DATA_DIR="/absolute/path/to/output" \
  bash scripts/run_batch.sh "/absolute/path/to/case.mac"
\end{lstlisting}

Replace the two paths with the desired output directory and macro file. The simulation messages appear in the terminal, and the output files declared by the macro are written to the specified directory.

Since the simulation overwrites existing files, using different target directories for different cases is recommended. If \path{OPNOVICE_DATA_DIR} is not set, the output is written to the \path{data/} directory in the root of the repository.

\section{Checking the run and using its output}

After the command finishes, check the terminal output for errors and for the line beginning \texttt{number of event = }. Its count should agree with the value requested by \path{/run/beamOn}. The output files are saved in the chosen directory. Which files are generated depends on the output selections and physics settings in the macro. Important CSV files include:
\begin{itemize}\setlength{\itemsep}{0pt}
  \item \path{results_nt_absorption.csv}: detected optical-photon positions;
  \item \path{results_nt_pr_int.csv}: primary-gamma interactions;
  \item \path{results_nt_scat_angles.csv}: Compton-scattering records;
  \item \path{results_nt_energy_deposition.csv}: energy-deposition steps in the scintillator.
\end{itemize}


\printbibliography[title={References}]

\end{document}
